\section{欧拉法}
\subsection{欧拉法}
考虑一阶常微分方程的初值问题\begin{equation}\label{equation:一阶常微分方程的初值问题.微分形式}
%@see: 《数值分析（第6版）》（李庆扬、王能超、易大义） P256 (9.1)
	\begin{cases}
		y' = f(x,y), \quad x \in [x_0,b], \\
		y(x_0) = y_0.
	\end{cases}
\end{equation}
如果我们用\(\increment y = f(x,y) \increment x\)近似代替\(\dd{y} = f(x,y) \dd{x}\)，
那么我们就可以用折线段\(\overline{P_0 \dotso P_n}\)
近似代替积分曲线\(y = y(x)\)，
其中\(P_0\)表示初始点\((x_0,y_0)\)，
其余点\(P_k(x_k,y_k)\ (k=1,2,\dotsc,n)\)满足\begin{equation*}
	\frac{y_k - y_{k-1}}{x_k - x_{k-1}}
	= f(x_{k-1},y_{k-1}),
\end{equation*}
或\begin{equation}
%@see: 《数值分析（第6版）》（李庆扬、王能超、易大义） P257 (9.4)
	y_k = y_{k-1} + h f(x_{k-1},y_{k-1}),
\end{equation}
并且\(x_1 - x_0 = x_2 - x_1 = \dotsb = x_n - x_{n-1} = h\).

为了分析计算过程中递推公式的精度，
假设准确值为\(y(x_k)\ (k=0,1,2,\dotsc,n)\)，
再进行泰勒展开得\begin{equation*}
	y(x_{k+1})
	= y(x_k + h)
	= y(x_k)
	+ h y'(x_k)
	+ \frac12 h^2 y''(\xi_k),
	\quad(x_k < \xi_k < x_{k+1}).
\end{equation*}
因为\(
	y'(x_k)
	= f(x_k,y(x_k))
\)，
所以误差为\begin{align*}
%@see: 《数值分析（第6版）》（李庆扬、王能超、易大义） P258 (9.6)
	y(x_{k+1}) - y_{k+1}
	&=
		y(x_k)
		+ h y'(x_k)
		+ \frac12 h^2 y''(\xi_k)
		- y_k
		- h f(x_k,y_k) \\
	&=
		[
			y(x_k)
			- y_k
		]
		+ h [
			y'(x_k)
			- f(x_k,y_k)
		]
		+ \frac12 h^2 y''(\xi_k),
\end{align*}
再假设\(x_k\)处的精确值与近似值恰好相等（即\(y(x_k) = y_k\)），
则\begin{equation}
	y(x_{k+1}) - y_{k+1}
	= \frac12 h^2 y''(\xi_k)
	\approx \frac12 h^2 y''(x_k).
\end{equation}
我们把这样算出的误差称为“欧拉法的\DefineConcept{局部截断误差}”.

\subsection{后退欧拉法}
利用牛顿--莱布尼茨公式可以从
一阶常微分方程 \labelcref{equation:一阶常微分方程的初值问题.微分形式}
得到\begin{equation}\label{equation:一阶常微分方程的初值问题.积分递推公式}
%@see: 《数值分析（第6版）》（李庆扬、王能超、易大义） P259 (9.7)
	y(x_{k+1})
	= y(x_k)
	+ \int_{x_k}^{x_{k+1}} f(t,y(t)) \dd{t}.
\end{equation}
显然欧拉法就相当于用左矩形公式\(h f(x_k,y_k)\)代替上式右端积分\(I_k \defeq \int_{x_k}^{x_{k+1}} f(t,y(t)) \dd{t}\).
类似地，如果用右矩形公式\(h f(x_{k+1},y_{k+1})\)代替\(I_k\)，
就能得到另一个公式\begin{equation}
%@see: 《数值分析（第6版）》（李庆扬、王能超、易大义） P259 (9.8)
	y_{k+1}
	= y_k
	+ h f(x_{k+1},y_{k+1})
	\quad(k=0,1,2,\dotsc).
\end{equation}
我们把这个递推公式称为\DefineConcept{后退欧拉法}.

后退欧拉法与欧拉法有着本质区别：
后者是关于\(y_{k+1}\)的一个直接的计算公式，
这类公式称为\DefineConcept{显式的}；
然而后退欧拉法的右端含有位置的\(y_{k+1}\)，
它实际上是关于\(y_{k+1}\)的一个函数方程，
这类公式称为\DefineConcept{隐式的}.
因此后退欧拉法也称为\DefineConcept{隐式欧拉法}.


显式方法与隐式方法各有特点.
有时考虑到数值稳定性等其他因素，人们需要选用隐式方法；
但是通常来说使用显式方法远比隐式方法方便.


后退欧拉法通常用迭代法求解，
其迭代过程的实质是逐步显式化.

假设我们首先用欧拉法算得\begin{equation*}
	y_{k+1}^{(0)}
	= y_k
	+ h f(x_k,y_k)
\end{equation*}
并将其作为迭代初值，
再代入后退欧拉法的递推公式中，使之转化为显式，
直接计算得\begin{equation*}
	y_{k+1}^{(1)}
	= y_k
	+ h f(x_{k+1},y_{k+1}^{(0)}),
\end{equation*}
然后继续将\(y_{k+1}^{(1)}\)代入后退欧拉法的递推公式中，
又有\begin{equation*}
	y_{k+1}^{(2)}
	= y_n
	+ h f(x_{k+1},y_{k+1}^{(1)}).
\end{equation*}
如此反复进行迭代，
得到递推数列\begin{equation}
%@see: 《数值分析（第6版）》（李庆扬、王能超、易大义） P259 (9.9)
	y_{k+1}^{(m+1)}
	= y_k
	+ h f(x_{k+1},y_{k+1}^{(m)})
	\quad(m=0,1,2,\dotsc).
\end{equation}
由于函数\(f\)关于\(y\)满足利普希茨条件，
那么上式与后退欧拉法递推公式相减得\begin{equation*}
	\abs{
		y_{k+1}^{(m+1)}
		- y_{k+1}
	}
	= h \abs{
		f(x_{k+1},y_{k+1}^{(m)})
		- f(x_{k+1},y_{k+1})
	}
	\leq h L
	\abs{
		y_{k+1}^{(m)}
		- y_{k+1}
	}.
\end{equation*}
由此可知，
只要取\(h\)使得\(h L < 1\)，
上述迭代显式化算法就能收敛到\(y_{k+1}\).

后退欧拉法的局部截断误差与欧拉法是相似的.

\subsection{梯形方法}
为了得到比欧拉法精度更高的计算公式，
我们可以用梯形公式近似代替积分\(I_k\)，
同时用\(y_k\)代替\(y(x_k)\)，
用\(y_{k+1}\)代替\(y(x_{k+1})\)，
便得\begin{equation*}
%@see: 《数值分析（第6版）》（李庆扬、王能超、易大义） P259 (9.10)
	y_{k+1}
	= y_k
	+ \frac{h}{2}
	[
		f(x_k,y_k)
		+ f(x_{k+1},y_{k+1})
	]
	\quad(k=0,1,2,\dotsc).
\end{equation*}
我们把上述递推公式称为\DefineConcept{梯形方法}.

梯形方法是隐式单步法，
可以用迭代法求解.
同后退欧拉法一样，
仍然用欧拉法提供迭代初值，
其迭代公式为\begin{equation}
%@see: 《数值分析（第6版）》（李庆扬、王能超、易大义） P260 (9.11)
	\begin{cases}
		y_{k+1}^{(0)}
		= y_k
		+ h f(x_k,y_k), \\
		y_{k+1}^{(m+1)}
		= y_k
		+ \frac{h}{2}
		[
			f(x_k,y_k)
			+ f(x_{k+1},y_{k+1}^{(m)})
		]
		\quad(k=0,1,2,\dotsc).
	\end{cases}
\end{equation}

为了分析迭代过程的收敛性，
将递推公式与迭代公式相减得\begin{equation*}
	y_{k+1}^{(m+1)}
	- y_{k+1}
	= \frac{h}{2}
	[
		f(x_{k+1},x_{k+1}^{(m)})
		- f(x_{k+1},y_{k+1})
	],
\end{equation*}
于是有\begin{equation*}
	\abs{
		y_{k+1}^{(m+1)}
		- y_{k+1}
	}
	\leq \frac{h L}{2}
	\abs{
		y_{k+1}^{(m)}
		- y_{k+1}
	}.
\end{equation*}
可以看出，只要取\(h\)使得\(h L / 2 < 1\)，
上述迭代显式化算法就能收敛到\(y_{k+1}\).

\subsection{改进欧拉法}
我们看到，梯形方法虽然提高了计算精度，
但是它的算法复杂.
在应用迭代公式进行实际计算时，每迭代一次，
都要重新计算函数\(f(x,y)\)的值，
而迭代又要反复进行若干次，计算量很大，
而且往往难以预测迭代何时停止.
为了控制计算量，通常只迭代一两次就转入下一步计算，这就简化了算法.

具体地说，我们先用欧拉法求出一个初步的近似值（称之为\DefineConcept{预测值}），
虽然其精度可能很差，
但是我们可以用梯形方法将它校正一次（将校正后算出的结果称为\DefineConcept{校正值}）：
\begin{equation}
%@see: 《数值分析（第6版）》（李庆扬、王能超、易大义） P260 (9.12)
	\begin{cases}
		% 预测
		\overline{y}_{k+1}
		= y_k
		+ h f(x_k,y_k), \\
		y_{k+1}
		= y_k
		+ \frac{h}{2}
		[
			f(x_k,y_k)
			+ f(x_{k+1},\overline{y}_{k+1})
		]
		\quad(k=0,1,2,\dotsc)
	\end{cases}
\end{equation}
或\begin{equation*}
	\begin{cases}
		y_p
		= y_k
		+ h f(x_k,y_k), \\
		y_c
		= y_k + h f(x_{k+1},y_p), \\
		y_{k+1}
		= \frac12 (y_p + y_c).
	\end{cases}
\end{equation*}
像这样建立的“预测--校正系统”通常称为\DefineConcept{改进欧拉法}.

\subsection{单步法的局部截断误差与阶}
一阶常微分方程的初值问题的单步法可以表示为如下一般形式：
\begin{equation}\label{equation:常微分方程初值问题的数值解法.单步法的一般形式}
%@see: 《数值分析（第6版）》（李庆扬、王能超、易大义） P261 (9.13)
	y_{k+1}
	= y_k
	+ h \phi(x_k,y_k,y_{k+1},h)
	\quad(k=0,1,2,\dotsc),
\end{equation}
其中多元函数\(\phi\)与\(f\)有关.
当\(\phi\)中含有\(y_{k+1}\)时，方法是隐式的；
当\(\phi\)中不含\(y_{k+1}\)时，方法是显式的.
因此，显式单步法可以表示为：
\begin{equation}\label{equation:常微分方程初值问题的数值解法.显式单步法的一般形式}
%@see: 《数值分析（第6版）》（李庆扬、王能超、易大义） P261 (9.14)
	y_{k+1}
	= y_k
	+ h \phi(x_k,y_k,h)
	\quad(k=0,1,2,\dotsc).
\end{equation}
我们把\(\phi(x,y,h)\)称为\DefineConcept{增量函数}.
例如，欧拉法的增量函数为\(\phi(x,y,h) = f(x,y)\).

\begin{definition}
%@see: 《数值分析（第6版）》（李庆扬、王能超、易大义） P261 定义9.1
设\(y\)是一阶常微分方程的初值问题的精确解，
把\begin{equation*}
	T_{k+1}
	\defeq
	y(x_{k+1})
	- y(x_k)
	- h \phi(x_k,y(x_k),h)
\end{equation*}
称为“显式单步法 \labelcref{equation:常微分方程初值问题的数值解法.显式单步法的一般形式} 的\DefineConcept{局部截断误差}”.
\end{definition}

误差\(T_{k+1}\)之所以被称为“局部的”，
是因为它隐含地假设\(x_k\)前的每一步都没有误差.
当第\(k\)步算得的近似值与精确值相等（即\(y_k = y(x_k)\)）时，只计算一步，就有\begin{equation*}
	y(x_{k+1}) - y_{k+1}
	= y(x_{k+1}) - [y_k + h \phi(x_k,y_k,h)]
	= y(x_{k+1}) - y(x_k) - h \phi(x_k,y(x_k),h)
	= T_{k+1}.
\end{equation*}
所以，局部截断误差可以理解为
使用显式单步法 \labelcref{equation:常微分方程初值问题的数值解法.显式单步法的一般形式}
计算一步的误差.

根据上式定义，欧拉法的局部截断误差为\begin{equation*}
	T_{k+1}
	= y(x_{k+1}) - y(x_k) - h f(x_k,y(x_k))
	= y(x_k + h)
	- y(x_k)
	- h y'(x_k)
	= \frac{h^2}{2} y''(x_k)
	+ O(h^3).
\end{equation*}
这里\(\frac{h^2}{2} y''(x_k)\)称为局部截断误差的主项.
显然\(T_{k+1} = O(h^2)\).

一般情形的局部截断误差主项的定义如下.
\begin{definition}
%@see: 《数值分析（第6版）》（李庆扬、王能超、易大义） P261 定义9.2
设\(y\)是一阶常微分方程的初值问题的精确解.
如果存在最大整数\(p\)
使得显式单步法 \labelcref{equation:常微分方程初值问题的数值解法.显式单步法的一般形式} 的
局部截断误差满足\begin{equation}\label{equation:常微分方程初值问题的数值解法.具有p阶精度的显式单步法需要满足的条件}
%@see: 《数值分析（第6版）》（李庆扬、王能超、易大义） P261 (9.15)
	T_{k+1}
	= O(h^{p+1}),
\end{equation}
则称“方法 \labelcref{equation:常微分方程初值问题的数值解法.显式单步法的一般形式}
具有\(p\)阶\DefineConcept{精度}”.
\end{definition}
%@see: 代数精度

如果微分方程的解是\(p\)次多项式，
那么具有\(p\)阶精度的方法能求出精确解.

若将\cref{equation:常微分方程初值问题的数值解法.具有p阶精度的显式单步法需要满足的条件}
展开，写成\begin{equation*}
	T_{k+1}
	= \psi(x_k,y(x_k)) h^{p+1}
	+ O(h^{p+2}),
\end{equation*}
则\(\psi(x_k,y(x_k)) h^{p+1}\)称为\DefineConcept{局部截断误差主项}.

\(p\)阶精度的定义
对隐式单步法也是适用的.
例如，后退欧拉法的局部截断误差为\begin{align*}
	T_{k+1}
	&= y(x_{k+1})
	- y(x_k)
	- h f(x_{k+1},y(x_{k+1})) \\
	&= h y'(x_k)
	+ \frac{h^2}{2} y''(x_k)
	+ O(h^3)
	- h [
		y'(x_k)
		+ h y''(x_k)
		+ O(h^2)
	] \\
	&= -\frac{h^2}{2} y''(x_k)
	+ O(h^3)
	= O(h^2).
\end{align*}
这里\(p=1\)，
说明隐式单步法具有1阶精度，其局部截断误差主项为\(-\frac{h^2}{2} y''(x_k)\).

同样地，对梯形方法有\begin{align*}
	T_{k+1}
	&= y(x_{k+1})
	- y(x_k)
	- h [y'(x_k) + y'(x_{k+1})] \\
	&= h y'(x_k)
	+ \frac{h^2}{2} y''(x_k)
	+ \frac{h^3}{3!} y'''(x_k)
	- \frac{h}{2}
	\left[
		y'(x_k)
		+ y'(x_k)
		+ h y''(x_k)
		+ \frac{h^2}{2} y'''(x_k)
	\right]
	+ O(h^4) \\
	&= -\frac{h^3}{12} y'''(x_k)
	+ O(h^4)
	= O(h^3).
\end{align*}
因此梯形方法具有2阶精度，
其局部误差主项为\(-\frac{h^3}{12} y'''(x_n)\).

后退欧拉法及梯形法这两种隐式方法都是从与微分方程等价的积分形式 \labelcref{equation:一阶常微分方程的初值问题.积分递推公式} 得到的.
由此可见，用数值积分中具有更高代数精度的求积公式（例如辛普森求积公式）
代替\cref{equation:一阶常微分方程的初值问题.积分递推公式} 右端积分，
虽然能使方法的精度提高，但由此建立的方法却不是显式的.
